% DERIVED from formal_layer.json — read-only, do not edit.
\begin{assumptionv}{ass:pomdp-kernel}[Joint reward and transition law]
The reward \(Y_t\) and next joint state \(S_{t+1}\) have conditional distribution
\[
\mathcal L(Y_t,S_{t+1}\mid\mathcal F_t^-,A_t)
=
K(\cdot\mid S_t,A_t),
\qquad t=1,\ldots,T,
\]
where \(\mathcal F_t^-\) is the full pre-action history information, \(A_t\) is the action, \(S_t\) is the current joint state, \(K\) is the joint reward and next-state law, and \(T\) is the observed trajectory length.
\end{assumptionv}

\begin{assumptionv}{ass:sequential-ignorability}[Observed-state sequential assignment]
Given the full pre-action history information \(\mathcal F_t^-\), the action \(A_t\) follows the behavior policy \(b\) at the observed state \(X_t\):
\[
\Pr(A_t=a\mid\mathcal F_t^-)=b(a\mid X_t),
\qquad a\in\mathcal A,\quad t=1,\ldots,T,
\]
where \(\mathcal A\) is the action space and \(T\) is the observed trajectory length.
\end{assumptionv}

\begin{assumptionv}{ass:stationary-start}[Stationary behavior-policy initial law]
The initial joint state \(S_1\) has the stationary joint-state law \(d_b\) under the behavior policy:
\[
S_1\sim d_b.
\]
\end{assumptionv}

\begin{assumptionv}{ass:policy-overlap}[Bounded action likelihood ratios]
The target action probabilities \(e(a\mid x)\) and behavior action probabilities \(b(a\mid x)\) satisfy
\[
e(a\mid x)\le Lb(a\mid x),
\qquad (a,x)\in\mathcal A\times\mathcal X,
\]
where \(L\) is the action likelihood-ratio bound, \(\mathcal A\) is the action space, and \(\mathcal X\) is the observed-state space.
\end{assumptionv}

\begin{assumptionv}{ass:behavior-contraction}[Behavior-policy transition contraction]
The behavior-policy joint-state transition matrix \(P_b\) satisfies
\[
\|\mu P_b-\mu'P_b\|_{\mathrm{TV}}
\le
\alpha\|\mu-\mu'\|_{\mathrm{TV}},
\qquad \mu,\mu'\in\Delta(\mathcal S),
\]
where \(\alpha\) is the contraction bound, \(\Delta(\mathcal S)\) is the probability simplex on the joint-state space \(\mathcal S\), and \(\|\cdot\|_{\mathrm{TV}}\) denotes total variation.
\end{assumptionv}

\begin{assumptionv}{ass:target-contraction}[Target-policy transition contraction]
The target-policy joint-state transition matrix \(P_e\) satisfies
\[
\|\mu P_e-\mu'P_e\|_{\mathrm{TV}}
\le
\alpha\|\mu-\mu'\|_{\mathrm{TV}},
\qquad \mu,\mu'\in\Delta(\mathcal S),
\]
where \(\alpha\) is the contraction bound, \(\Delta(\mathcal S)\) is the probability simplex on the joint-state space \(\mathcal S\), and \(\|\cdot\|_{\mathrm{TV}}\) denotes total variation.
\end{assumptionv}

\begin{definitionv}{def:binary-pomdp-class}[Binary POMDP experiment class]
The class \(\mathcal M_T^{(2)}(t_0,\zeta)\) consists of binary trajectory experiments of observed length \(T\in\mathbb N\), with observed-state, hidden-state, and action alphabets
\[
\mathcal X=\mathcal H=\mathcal A=\{0,1\},
\qquad \mathcal Y=[0,1].
\]
Each experiment comprises a joint reward and next-state law \(K\), behavior and target action probabilities \(b\) and \(e\), and a full trajectory probability law inducing the law of \(\mathsf O_T\), the observed state, action, and reward trajectory. Membership is defined by the following conditions:
\begin{itemize}
\item \textbf{(Parameter range.)} The mixing-scale parameter \(t_0\) and logarithmic action likelihood-ratio bound \(\zeta\) satisfy
\[
t_0>0,\qquad \zeta\ge 0.
\]
The associated contraction and overlap bounds are
\[
\alpha=\exp(-1/t_0),\qquad L=\exp(\zeta).
\]
\item \textbf{(Joint law.)} The experiment satisfies \cref{obj:ass:pomdp-kernel}.
\item \textbf{(Sequential assignment.)} The experiment satisfies \cref{obj:ass:sequential-ignorability}.
\item \textbf{(Stationary start.)} The experiment satisfies \cref{obj:ass:stationary-start}.
\item \textbf{(Policy overlap.)} The policies satisfy \cref{obj:ass:policy-overlap} with overlap bound \(L\).
\item \textbf{(Behavior contraction.)} The behavior-induced joint-state transition matrix satisfies \cref{obj:ass:behavior-contraction} with contraction bound \(\alpha\).
\item \textbf{(Target contraction.)} The target-induced joint-state transition matrix satisfies \cref{obj:ass:target-contraction} with contraction bound \(\alpha\).
\end{itemize}
\end{definitionv}

\begin{definitionv}{def:target-value}[Stationary target value $\theta(K,e)$]
The stationary target mean reward is
\[
\theta(K,e)
:=
d_e r_e
=
\sum_{s=(x,h)\in\mathcal S}
d_e(s)
\sum_{a\in\mathcal A}
e(a\mid x)
\sum_{s'\in\mathcal S}
\int_{\mathcal Y}
y\,K(dy,s'\mid s,a),
\]
where \(d_e\) is the stationary joint-state law under the target policy, \(r_e\) is the vector of target-policy conditional mean rewards, \(\mathcal S\) is the four-point joint-state space, \(e(a\mid x)\) is the target action probability at observed state \(x\), and \(K(dy,s'\mid s,a)\) is the joint reward and next-state law.
\end{definitionv}

\begin{definitionv}{def:observable-moments}[Observable moments $m_k$ and $\widehat m_k$]
The weighted reward scores, their population moments, and their common-window empirical moments are, for \(k=0,\ldots,6\),
\[
Z_t(k)
:=
Y_t\prod_{j=t-k}^{t}\rho_j,
\qquad
m_k
:=
\mathbb E_b[Z_t(k)],
\qquad
\widehat m_k
:=
\frac{1}{T-6}
\sum_{t=7}^{T}Z_t(k).
\]
Here \(Y_t\) is the bounded reward at time \(t\), \(\rho_j\) is the observed action likelihood ratio at time \(j\), \(\mathbb E_b\) denotes expectation under the stationary behavior experiment, and \(T\) is the observed trajectory length.
\end{definitionv}

\begin{algorithmv}{def:stable-grid-estimator}[Stable grid estimator $\widehat\theta_T$]
The inputs are the empirical moments \(\widehat m_0,\ldots,\widehat m_6\), the trajectory length \(T\), and the contraction bound \(\alpha\).
\begin{enumerate}
\item Set the grid resolution, probability grid, uniform transition matrix, and stabilization weight to
\[
M:=\left\lceil(22+36/\alpha)T\right\rceil,
\qquad
\mathcal G_M
:=
\left\{
u\in[0,1]^4:
\sum_{i=1}^{4}u_i=1,\ 
u_i\in M^{-1}\mathbb Z\ \text{for all }i
\right\},
\]
\[
U:=\frac{\mathbf1\mathbf1^{\mathsf T}}{4},
\qquad
\lambda:=\frac{\alpha}{\alpha+6/M},
\]
where \(\mathbf1\) is the four-dimensional vector of ones.
\item Enumerate the candidate triples in
\[
\mathcal C_M
:=
\left\{
(\nu,R,r):
\begin{gathered}
\nu\in\mathcal G_M,\quad
R_{i,\cdot}\in\mathcal G_M\ \text{for }i=1,\ldots,4,\\
\delta(R)\le\alpha+6/M,\quad
r\in\{0,1/M,\ldots,1\}^{4}
\end{gathered}
\right\},
\]
where \(\delta(R)\) is the Dobrushin contraction coefficient. For each candidate, form
\[
P(R):=\lambda R+(1-\lambda)U.
\]
\item Select the lexicographically first triple \((\nu_*,R_*,r_*)\) among the minimizers
\[
(\nu_*,R_*,r_*)
\in
\operatorname*{arg\,min}_{(\nu,R,r)\in\mathcal C_M}
\max_{0\le k\le6}
\left|\widehat m_k-\nu P(R)^k r\right|.
\]
\item Form the selected transition matrix and output its stationary reward:
\[
P_*:=P(R_*),
\qquad
\widehat\theta_T:=\pi(P_*)r_*,
\]
where \(\pi(P_*)\) is the stationary probability vector satisfying
\[
\pi(P_*)P_*=\pi(P_*).
\]
\end{enumerate}
\end{algorithmv}

\begin{definitionv}{def:minimax-risk}[Minimax risk $R_T^{(2)}(t_0,\zeta)$]
The minimax mean-squared error is
\[
R_T^{(2)}(t_0,\zeta)
:=
\inf_{\substack{\widetilde\theta\colon \mathsf O_T\to[0,1]\ \mathrm{measurable}\\
\mathrm{with\ respect\ to}\ \sigma(\mathsf O_T,b,e,t_0,\zeta)}}
\sup_{M\in\mathcal M_T^{(2)}(t_0,\zeta)}
\mathbb E_M\!\left[(\widetilde\theta-\theta)^2\right].
\]
The infimum ranges over measurable observed-data procedures taking values in \( [0,1] \). Here \(\mathsf O_T\) is the observed trajectory of length \(T\); \(b\) and \(e\) are the supplied behavior and target policies; \(t_0\) and \(\zeta\) are the supplied mixing scale and logarithmic action likelihood-ratio bound; \(\mathcal M_T^{(2)}(t_0,\zeta)\) is the binary experiment class; \(\mathbb E_M\) denotes expectation under experiment \(M\); and \(\theta\) is that experiment's stationary target value.
\end{definitionv}

\begin{theoremv}{thm:stable-grid-upper}[Uniform stable-grid risk bound]
Let the mixing-scale parameter \(t_0\), logarithmic action likelihood-ratio bound \(\zeta\), and trajectory length \(T\) satisfy:
\begin{itemize}
\item \textbf{(Mixing scale.)} \(t_0\in\mathbb R\) and \(t_0>0\).
\item \textbf{(Overlap bound.)} \(\zeta\in\mathbb R\) and \(\zeta\geq0\).
\item \textbf{(Horizon.)} \(T\in\mathbb N\) and \(T\geq12\).
\end{itemize}
Define the contraction bound \(\alpha\), action likelihood-ratio bound \(L\), stationary-value modulus constant \(B_\alpha\), and moment variance constant \(V_{\alpha,L}\) by
\[
\alpha=\exp(-1/t_0),\qquad
L=\exp(\zeta),\qquad
B_\alpha=\left(\frac{1+\alpha}{1-\alpha}\right)^6,\qquad
V_{\alpha,L}=13L^7+\frac{2}{1-\alpha}.
\]
For every binary-state, binary-action experiment \(M\in\mathcal M_T^{(2)}(t_0,\zeta)\), the stable-grid estimator \(\widehat\theta_T\) of \cref{obj:def:stable-grid-estimator}, using that experiment's behavior and target policies, and the stationary target value \(\theta\) of \cref{obj:def:target-value} satisfy
\[
\mathbb E_M\!\left[(\widehat\theta_T-\theta)^2\right]
\leq
\frac{B_\alpha^2\bigl(112V_{\alpha,L}+2\bigr)}{T},
\]
where \(\mathbb E_M\) denotes expectation under the experiment's observed trajectory law.
\end{theoremv}

\begin{definitionv}{def:four-state-testing-family}[Four-state testing family]
The four-state testing family is indexed by the trajectory length \(T\in\{0,1,\ldots\}\) and a signed perturbation \(v\in\{-v_T,+v_T\}\) satisfying \(|v|\le 1/16\), where the testing perturbation amplitude is
\[
v_T=
\begin{cases}
0,&T=0,\\
\dfrac{1}{16\sqrt T},&T\ge 1.
\end{cases}
\]
The observed state, hidden state, and action take values in \(\{0,1\}\), so the joint-state space is \(\mathcal S=\{0,1\}^2\). The behavior and target policies are
\[
b(a\mid x)=e(a\mid x)=\frac12,
\qquad x,a\in\{0,1\}.
\]
The observation probabilities, reward probabilities, and successor hidden-state probabilities are defined by
\[
E(x\mid h)=
\begin{cases}
3/5,&x=h,\\
2/5,&x\ne h,
\end{cases}
\qquad
q_v(h)=\frac12+\frac{2h-1}{8}+v,
\qquad
p(a)=\frac14+\frac a2.
\]
The joint reward and next-state kernel has mass
\[
K_v^{(T)}(y,x',h'\mid x,h,a)
=
\operatorname{Ber}_{q_v(h)}(y)
\operatorname{Ber}_{p(a)}(h')
E(x'\mid h'),
\qquad y,x',h'\in\{0,1\},
\]
where the Bernoulli masses are
\[
\operatorname{Ber}_{q_v(h)}(y)=
\begin{cases}
q_v(h),&y=1,\\
1-q_v(h),&y=0,
\end{cases}
\qquad
\operatorname{Ber}_{p(a)}(h')=
\begin{cases}
p(a),&h'=1,\\
1-p(a),&h'=0.
\end{cases}
\]
The initial joint state \(S_1\) has distribution
\[
\Pr\{S_1=(x,h)\}=\frac12 E(x\mid h).
\]
The full trajectory law is the finite mixture of Dirac masses over binary state, action, and reward paths, with each path weighted by its initial joint-state mass multiplied by its behavior-policy action probabilities and joint reward and next-state kernel masses.
\end{definitionv}

\begin{theoremv}{thm:fixed-binary-minimax}[Fixed binary minimax rate]*
Let \(t_0\) be the mixing-scale parameter, \(\zeta\) the logarithmic action likelihood-ratio bound, and \(T\) the observed trajectory length, with
\begin{itemize}
\item \textbf{(Positive mixing scale.)} \(t_0>0\).
\item \textbf{(Nonnegative overlap exponent.)} \(\zeta\geq0\).
\item \textbf{(Horizon range.)} \(T\in\mathbb N\) and \(T\geq12\).
\end{itemize}
Define the contraction and action-overlap bounds, and the associated modulus and variance constants, by
\[
\alpha=\exp(-1/t_0),\qquad L=\exp(\zeta),\qquad
B_\alpha=\left(\frac{1+\alpha}{1-\alpha}\right)^6,\qquad
V_{\alpha,L}=13L^7+\frac{2}{1-\alpha}.
\]
The observable-data minimax mean-squared risk in \cref{obj:def:minimax-risk} satisfies
\[
\frac{1}{1024T}
\leq R_T^{(2)}(t_0,\zeta)
\leq \frac{B_\alpha^2(112V_{\alpha,L}+2)}{T}.
\]

Define the testing perturbation amplitude by
\[
v_T=\frac{1}{16\sqrt T}.
\]
Both experiments in \cref{obj:def:four-state-testing-family} with perturbations \(v=+v_T\) and \(v=-v_T\) belong to the binary class \(\mathcal M_T^{(2)}(t_0,\zeta)\) of \cref{obj:def:binary-pomdp-class}. In each experiment the behavior and target policies are equal fair policies, and their stationary joint-state laws coincide:
\[
e(a\mid x)=b(a\mid x)=\frac12,\qquad d_e=d_b.
\]
For every stationary occupancy-ratio bound \(C>1\), both experiments satisfy
\[
d_e(s)\leq C\,d_b(s)
\qquad\text{for every joint state }s.
\]
For every observable-data estimator \(\widetilde\theta\) taking values in \([0,1]\), the maximum of its mean-squared risks over these two experiments is at least \(1/(1024T)\).

If additionally \(\zeta>0\), then for every \(C>1\), define the cardinality-uniform rate exponent and stationary-overlap radius by
\[
\beta=\frac{2}{2+t_0\zeta},\qquad q_C=\frac{C-1}{C}.
\]
Then
\[
0<\beta<1,\qquad
\lim_{\substack{n\to\infty\\ n\in\mathbb N}}
\frac{n^{-1}}{n^{-\beta}q_C^{\,2(1-\beta)}}=0.
\]
There exist a positive reset lower-bound constant \(c_{\mathrm{Reset}}>0\) and a depth schedule \(Q:\mathbb N\to\mathbb N\), together with constants \(a_1,a_2>0\), such that, for every sufficiently large \(n\),
\[
a_1\log n\leq Q(n)\leq a_2\log n.
\]
For every sufficiently large \(n\), one has \(n\geq1\) and \(Q(n)\geq1\), and the two signed depth-\(Q(n)\) reset experiments of \citet{CausalSmithLatentOverlap2026} each have exactly \(2(Q(n)+1)\) hidden states. For every observable-data estimator taking values in \([-1,1]\) in that cardinality-uniform decision problem, the maximum of its mean-squared risks over these two reset experiments is at least
\[
c_{\mathrm{Reset}}\,n^{-\beta}.
\]
For every sufficiently large \(n\), the negative-sign reset experiment has hidden-state cardinality
\[
2\bigl(Q(n)+1\bigr)>2.
\]
\end{theoremv}
% CAUSALSMITH-CITED-SCOPE-BEGIN thm:fixed-binary-minimax
\verificationfootnotetext{\textbf{Formalization scope.} This result uses the published conclusion \citep{CausalSmithLatentOverlap2026} (Main minimax theorem and reset-chain lower-bound construction, as supplied in the frozen target and Topic record.). The cited source proof is not formalized here: Lean verifies its use and all remaining steps. If that cited conclusion is false, the portions of this result that depend on it are not certified.}
% CAUSALSMITH-CITED-SCOPE-END thm:fixed-binary-minimax

\begin{propositionv}{prop:coincident-policy-reduction}[Coincident policy reduction]
Consider a trajectory experiment with binary observed state, binary hidden state, and binary action. Suppose:
\begin{itemize}
\item \textbf{(Positive mixing scale.)} The mixing-scale parameter \(t_0\) satisfies \(t_0>0\), and the associated contraction factor is
\[
\alpha=\exp(-1/t_0).
\]
\item \textbf{(Nonempty horizon.)} The trajectory length \(T\) is an integer with \(T\geq 1\).
\item \textbf{(Kernel law.)} \cref{obj:ass:pomdp-kernel} holds.
\item \textbf{(Sequential assignment.)} \cref{obj:ass:sequential-ignorability} holds.
\item \textbf{(Stationary start.)} \cref{obj:ass:stationary-start} holds.
\item \textbf{(Behavior contraction.)} \cref{obj:ass:behavior-contraction} holds with contraction factor \(\alpha\).
\item \textbf{(Coincident policies.)} The target and behavior policies satisfy \(e=b\).
\end{itemize}
The induced joint-state transition matrices coincide. Explicitly, for all joint states \(s,s'\in\mathcal S\), the four-point joint-state space,
\[
P_e(s,s')
=
\sum_{a\in\{0,1\}}e(a\mid x)
   \int_{[0,1]}K(dy,s'\mid s,a)
=
\sum_{a\in\{0,1\}}b(a\mid x)
   \int_{[0,1]}K(dy,s'\mid s,a)
=
P_b(s,s'),
\]
where \(x\) is the observed coordinate of \(s\). Their stationary joint-state laws, \(d_e\) and \(d_b\), also coincide:
\[
d_e=d_b.
\]
Writing \(\theta\) for the stationary target value in \cref{obj:def:target-value}, the depth-zero population weighted-reward moment \(m_0\) satisfies, at every observed epoch,
\[
m_0=\mathbb E_b[Y_t]=\theta,
\qquad t=1,\ldots,T,
\]
where \(Y_t\) is the observed reward and \(\mathbb E_b\) denotes expectation under the behavior experiment. The sample mean is unbiased and obeys
\[
\mathbb E_b\!\left[T^{-1}\sum_{t=1}^{T}Y_t\right]=\theta,
\qquad
\mathbb E_b\!\left[
\left\{T^{-1}\sum_{t=1}^{T}Y_t-\theta\right\}^{2}
\right]
\leq
\frac{1+2/(1-\alpha)}{T}.
\]
\end{propositionv}

\begin{propositionv}{prop:exhaustive-grid-arithmetic-complexity}[Exact grid candidate counts]
Fix \(t_0>0\), the mixing-scale parameter, and define the associated contraction bound by
\[
\alpha=\exp(-1/t_0).
\]
For every integer \(M\ge0\), the four-coordinate simplex grid, represented by integer coordinates in \(\{0,\ldots,M\}\) summing to \(M\), has cardinality
\[
\binom{M+3}{3}.
\]
The unfiltered candidate list consists of an initial simplex vector \(\nu\), four simplex rows forming the transition candidate \(R\), and a reward vector \(r\) with four independently chosen grid coordinates. Its cardinality is exactly
\[
\binom{M+3}{3}^{5}(M+1)^4.
\]
For every integer \(M\ge0\) and every real contraction threshold \(\alpha\), retaining the candidates satisfying
\[
\delta(R)\le \alpha+\frac{6}{M},
\]
where \(\delta(R)\) is the Dobrushin contraction coefficient and division by zero is interpreted as zero, yields a cardinality at most
\[
\binom{M+3}{3}^{5}(M+1)^4.
\]

There exists a constant \(c>0\) such that, for every integer trajectory length \(T\ge1\), at the resolution
\[
M=\left\lceil(22+36/\alpha)T\right\rceil,
\qquad \alpha=\exp(-1/t_0),
\]
the unfiltered candidate-list cardinality and the filtered candidate-list cardinality are each at most \(cT^{19}\).
\end{propositionv}

\begin{definitionv}{def:polynomial-time-handle}[Stable fitting problem \(\mathfrak H_{\mathrm{poly}}\)]
The problem \(\mathfrak H_{\mathrm{poly}}\) is to construct an algorithm with an explicit cost bound in a specified arithmetic or bit-complexity model, including precision requirements, satisfying the following conditions:
\begin{itemize}
\item \textbf{(Input.)} The algorithm starts from the \(3\times3\) Hankel matrix
\[
\bigl(\widehat m_{i+j+1}-\widehat m_{i+j}\bigr)_{i,j=0}^{2},
\]
where \(\widehat m_k\) denotes the common-window empirical weighted reward moment.
\item \textbf{(Rank adaptation.)} It adapts over recurrence ranks zero through three.
\item \textbf{(Stability.)} It enforces transient roots in
\[
\{z:|z|\leq\alpha\},
\]
where \(\alpha\) is the contraction bound.
\item \textbf{(Objective accuracy.)} It returns a stable four-state scalar realization whose seven-moment objective is within \(T^{-1}\) of the global optimum, with a guarantee of this gap, where \(T\) is the observed trajectory length.
\item \textbf{(Uniformity.)} The guarantee is uniform over rank-deficient strata, with positive singular-value margins allowed to vanish.
\end{itemize}
For fixed \(t_0\), the mixing-scale parameter, the exhaustive estimator is specified by a finite candidate list of size \(O(T^{19})\). Its stated computational bound concerns this candidate count; runtime and arithmetic-operation bounds for fitting and selection form part of the computational target.
\end{definitionv}

\begin{remarkv}{oeq:polynomial-time-attainment}[Rank-adaptive computational attainment]
A natural computational question is whether \(\mathfrak H_{\mathrm{poly}}\), the cost-bounded stable fitting problem, can implement fitting and lexicographic selection over, or replace, the \(O(T^{19})\)-sized stable-grid candidate list by a rank-adaptive algorithm with explicit runtime, bit-complexity, and precision semantics. The required output is a stable four-state scalar realization whose seven-moment objective is guaranteed to be within \(T^{-1}\) of the global optimum, uniformly over \(\mathcal M_T^{(2)}(t_0,\zeta)\), the class of binary-state and binary-action experiments satisfying the six law conditions. Here \(T\) is the observed trajectory length, \(t_0\) is the mixing-scale parameter, and \(\zeta\) is the logarithmic action likelihood-ratio bound. The question requires uniformity over every rank-deficient Hankel stratum, with positive singular-value margins allowed to vanish.
\end{remarkv}

\begin{lemmav}{lem:pair-polynomial-modulus}[Pair-polynomial identity and value modulus]
Let \(\mathcal S=\{0,1\}\times\{0,1\}\) be the four-state joint-state space. Let \(P,P'\) be real matrices indexed by \(\mathcal S\), and let \(\nu,\nu',\pi,\pi',r,r'\) be real row vectors indexed by \(\mathcal S\). Suppose:
\begin{itemize}
\item \textbf{(Contraction parameter.)} The contraction bound \(\alpha\) satisfies \(0<\alpha<1\).
\item \textbf{(Transition matrices.)} Both \(P\) and \(P'\) are row-stochastic: their entries are nonnegative and each row sums to one.
\item \textbf{(Initial distributions.)} Both \(\nu\) and \(\nu'\) are probability vectors.
\item \textbf{(Stationary distributions.)} Both \(\pi\) and \(\pi'\) are probability vectors, with
\[
\pi P=\pi,\qquad \pi'P'=\pi'.
\]
\item \textbf{(Rewards.)} For every \(s\in\mathcal S\),
\[
0\le r(s)\le1,\qquad 0\le r'(s)\le1.
\]
\item \textbf{(Contraction bounds.)} The Dobrushin coefficients satisfy
\[
\begin{aligned}
\delta(P)&=\max_{i,j\in\mathcal S}\frac12
 \sum_{s\in\mathcal S}|P_{is}-P_{js}|\le\alpha,\\
\delta(P')&=\max_{i,j\in\mathcal S}\frac12
 \sum_{s\in\mathcal S}|P'_{is}-P'_{js}|\le\alpha.
\end{aligned}
\]
\end{itemize}
Then eigenvalue \(1\) has algebraic multiplicity one for each of \(P\) and \(P'\). Define the pair transient characteristic polynomial and its coefficients by
\[
Q_{P,P'}(z)
=\frac{\det(z\,\mathrm{Id}-P)}{z-1}
 \frac{\det(z\,\mathrm{Id}-P')}{z-1}
=\sum_{k=0}^{6}Q_k z^k,
\]
where \(\mathrm{Id}\) is the identity matrix on \(\mathcal S\), and define the modulus constant by
\[
B_\alpha=\left(\frac{1+\alpha}{1-\alpha}\right)^6.
\]
With products interpreted as
\[
\pi r=\sum_{s\in\mathcal S}\pi(s)r(s),
\qquad
\nu P^k r=\sum_{s\in\mathcal S}(\nu P^k)(s)r(s),
\]
and analogously for the primed vectors, the exact identity and bound are
\[
Q_{P,P'}(1)(\pi r-\pi'r')
=\sum_{k=0}^{6}Q_k\bigl(\nu P^k r-\nu'P'^k r'\bigr)
\]
and
\[
|\pi r-\pi'r'|
\le B_\alpha
 \max_{0\le k\le6}\bigl|\nu P^k r-\nu'P'^k r'\bigr|.
\]
\end{lemmav}

\begin{lemmav}{lem:observable-intervention-moments}[Observable intervention moment bounds]
Consider an experiment in \(\mathcal M_T^{(2)}(t_0,\zeta)\), the binary experiment class of \cref{obj:def:binary-pomdp-class}, with trajectory length \(T\). Set
\[
L=\exp(\zeta),\qquad \alpha=\exp(-1/t_0).
\]
On the four-point joint-state space \(\mathcal S\), let \(P_b\) and \(P_e\) be the behavior- and target-policy transition matrices, respectively, and let \(r_e\) be the target-policy conditional mean reward vector:
\[
\begin{aligned}
P_b(s,s')&=\sum_{a\in\{0,1\}}b(a\mid x)
                  \int_{\mathcal Y}K(dy,s'\mid s,a),\\
P_e(s,s')&=\sum_{a\in\{0,1\}}e(a\mid x)
                  \int_{\mathcal Y}K(dy,s'\mid s,a),\\
r_e(s)&=\sum_{a\in\{0,1\}}e(a\mid x)
                  \sum_{s'\in\mathcal S}\int_{\mathcal Y}y\,K(dy,s'\mid s,a),
\qquad s=(x,h).
\end{aligned}
\]
Here \(K\) is the joint reward and successor-state law, \(\mathcal Y=[0,1]\) is the reward space, and \(d_b\) is the stationary probability row vector of \(P_b\).

Use zero-based epochs \(0,\ldots,T-1\). For observed state \(X_j\), action \(A_j\), and reward \(Y_j\), define the observed likelihood ratio and weighted reward score by
\[
\rho_j=
\begin{cases}
e(A_j\mid X_j)/b(A_j\mid X_j),& b(A_j\mid X_j)>0,\\
0,& b(A_j\mid X_j)=0,
\end{cases}
\qquad
Z_t(k)=Y_t\prod_{j=t-k}^{t}\rho_j
\quad (k\le t).
\]
Write \(\mathbb E_b\) for expectation under the experiment's stationary behavior law. For every \(k\in\{0,\ldots,6\}\) and every epoch \(t\) with \(k\le t<T\),
\[
m_k=\mathbb E_b[Z_t(k)]
    =d_bP_e^kr_e
    =\sum_{s\in\mathcal S}(d_bP_e^k)(s)\,r_e(s),
\qquad
\mathbb E_b[Z_t(k)^2]\le L^{k+1}.
\]
Moreover, for every integer \(h\ge k+1\) such that \(t+h<T\),
\[
\left|
\mathbb E_b\!\left[
\bigl(Z_t(k)-\mathbb E_b[Z_t(k)]\bigr)
\bigl(Z_{t+h}(k)-\mathbb E_b[Z_{t+h}(k)]\bigr)
\right]
\right|
\le \alpha^{h-k-1}.
\]
\end{lemmav}

\begin{lemmav}{lem:testing-family-membership}[Testing family membership and separation]
Let \(T\) be the observed trajectory length, \(t_0\) the mixing-scale parameter, and \(\zeta\) the logarithmic action likelihood-ratio bound, with
\begin{itemize}
\item \textbf{(Length.)} \(T\in\mathbb N\) and \(T\geq12\).
\item \textbf{(Mixing scale.)} \(t_0>0\).
\item \textbf{(Overlap.)} \(\zeta\geq0\).
\end{itemize}
Define the testing perturbation amplitude by
\[
v_T=\frac{1}{16\sqrt T}.
\]
For each \(v\in\{-v_T,+v_T\}\), the experiment of \cref{obj:def:four-state-testing-family}, with kernel \(K_v^{(T)}\), the specified initial and trajectory laws, and fair coincident policies
\[
e(a\mid x)=b(a\mid x)=\frac12,
\]
belongs to \(\mathcal M_T^{(2)}(t_0,\zeta)\), the class satisfying the six law conditions in \cref{obj:def:binary-pomdp-class}. Its stationary target values, defined in \cref{obj:def:target-value}, satisfy
\[
\left|\theta(K_{+v_T}^{(T)},e)-\theta(K_{-v_T}^{(T)},e)\right|
=\frac{1}{8\sqrt T}.
\]
Writing \(\mathbb P_{+}^{\mathsf O_T}\) and \(\mathbb P_{-}^{\mathsf O_T}\) for the respective laws projected onto the observed state, action, and reward trajectory, their relative entropy is finite and satisfies
\[
\operatorname{KL}\!\left(
\mathbb P_{+}^{\mathsf O_T}\Vert\mathbb P_{-}^{\mathsf O_T}
\right)\leq\frac1{12}.
\]
For each of the two experiments, the stationary joint-state laws \(d_e\) and \(d_b\) under the target and behavior policies coincide:
\[
d_e=d_b.
\]
Moreover, for every stationary occupancy-ratio constant \(C>1\), both experiments satisfy
\[
d_e(s)\leq C\,d_b(s)
\qquad\text{for every }s\in\mathcal S,
\]
where \(\mathcal S\) is the four-point joint-state space.
\end{lemmav}
